\subsection{Dissipation and mobility choices}
\label{sec:dissipation-mobilities}

\subsubsection{Modified energy law}

Define, for every \(n\),
\begin{equation}
  \Hcal_{\mathrm{mom}}^n
  =(q^n)^2+\frac{1}{2\eta}\norm{p^n}_{M^n}^2.
  \label{eq:abstract-momentum-energy}
\end{equation}

\begin{theorem}
\label{thm:monotone-mobility-dissipation}
For \(0\leq\beta<1\) and \(0\prec M^{n+1}\preceq M^n\), the update
\cref{eq:covector-map,eq:abstract-momentum-step,eq:abstract-momentum-states}
satisfies
\begin{equation}
\begin{aligned}
  \Hcal_{\mathrm{mom}}^n
  -\Hcal_{\mathrm{mom}}^{n+1}
  ={}&\frac{\norm{p^n}_{M^n-M^{n+1}}^2}{2\eta}
  +\frac{\norm{p^{n+1}-\beta p^n}_{M^{n+1}}^2}{2\eta}\\
  &+\frac{1-\beta^2}{2\eta}\norm{p^n}_{M^{n+1}}^2
  +(d^{n+1})^2
  +\alpha\norm{\Delta\theta^{n+1}}_{S_m^n}^2
  \geq0.
\end{aligned}
\label{eq:monotone-mobility-dissipation}
\end{equation}
\end{theorem}

The proof is a weighted polarization of
\cref{eq:abstract-momentum-balance} against
\(\Delta\theta^{n+1}=M^{n+1}p^{n+1}\), followed by substitution of
\cref{eq:common-identities}; it is given in \cref{app:energy-law-proof}.

\begin{remark}
\label{rem:dissipated-quantity}
The dissipated quantity is the modified energy \(\Hcal_{\mathrm{mom}}^n\),
which controls the scalar and kinetic energies together rather than the
objective \(F(\theta^n)\) alone. Its five nonnegative terms are, in order,
the dissipation from the mobility change, the inertial residual, the momentum
damping, the scalar tracking term, and the curvature gap, which the
relaxation step of \cref{sec:relaxation-convergence} sums.
\end{remark}

\subsubsection{Euclidean momentum}

Set \(M^n=I_d\) and identify \(p^n=s^n\).  The abstract state equation becomes
\begin{equation}
  \left(I_d+\eta B_\alpha^n\right)s^{n+1}
  =\beta s^n-\eta\frac{q^n}{Q^n}g^n.
  \label{eq:euclidean-momentum-step}
\end{equation}
The constant mobility \(M^n\equiv I_d\) satisfies
\cref{eq:abstract-mobility-order} trivially, so
\cref{thm:monotone-mobility-dissipation} applies for every
\(0\leq\beta<1\); its mobility-change term vanishes, and the modified
energy is \((q^n)^2+\norm{s^n}^2/(2\eta)\).

\subsubsection{AMSGrad-type diagonal mobility}

An AMSGrad-type maximum envelope \cite{reddi2018amsgrad} gives a diagonal
mobility that is nonincreasing in the Loewner order. Let \(0\leq\beta_2<1\),
\(\epsilon>0\), and initialize
\(v^0=\overline v^0=0\).  With \(\odotv\) the componentwise product and the
square root, division, and maximum below taken componentwise, define
\begin{subequations}
\label{eq:adaptive-second-moment}
\begin{align}
  v^{n+1}
  &=\beta_2v^n+(1-\beta_2)(g^n\odotv g^n),
  \label{eq:adaptive-v}\\
  \widehat v^{n+1}
  &=\frac{v^{n+1}}{1-\beta_2^{n+1}},
  \label{eq:adaptive-bias-correction}\\
  \overline v^{n+1}
  &=\max\{\overline v^n,\widehat v^{n+1}\},
  \label{eq:adaptive-max}\\
  M^{n+1}
  &=\diag\left(
  \frac{1}{\sqrt{\overline v^{n+1}}+\epsilon}
  \right).
  \label{eq:adaptive-mobility}
\end{align}
\end{subequations}
Then \(M^0=\epsilon^{-1}I_d\), and the maximum envelope implies
\begin{equation}
  0\prec M^{n+1}\preceq M^n,
  \label{eq:adaptive-mobility-order}
\end{equation}
so \cref{thm:monotone-mobility-dissipation} applies to this mobility for
every \(0\leq\beta<1\).
